\begin{abstract}
The rapid growth of mobile application technology has created new opportunities for delivering psychological assessments in a more accessible and efficient manner compared to conventional methods; however, a major challenge lies in translating complex interface designs into mobile applications that are consistent, responsive, and user-friendly. The Jatidiri application is developed as a mobile-based platform that provides various psychological tests, including intelligence, self-control, talent, happiness, anxiety, learning style, stress management, psychological health, potential, career, adult intelligence, and resilience tests, to support users in understanding their psychological condition and personal potential. This study aims to design and implement the user interface of the Jatidiri mobile application using the Flutter framework based on UI designs created in Figma, by applying the Prototype development method to enable iterative and adaptive development.

The Prototype method is applied through the stages of requirement gathering, prototype development, prototype evaluation and refinement, and final implementation and testing. The development process focuses on translating the Figma-based UI designs into Flutter components while maintaining layout consistency, color schemes, typography, and user interaction patterns. Black box testing is conducted to ensure that each feature and interface element functions according to the specified requirements and design.

The results show that the Prototype method is effective in supporting the implementation of Figma designs into a Flutter-based mobile application, resulting in a responsive and visually consistent user interface. However, this study is limited in scope, as it does not include psychometric validation of the test content nor formal usability testing with end users; therefore, the findings primarily reflect the technical and design implementation aspects of the application.
\end{abstract}

\noindent \textbf{Kata kunci:} Aplikasi Mobile; Tes Psikologi; Flutter; Metode Prototype; UI/UX.
